% `tex.web` §294 `print_meaning`, and the one `print_ln` that tells `\show`
% from `\meaning`:
%
%     procedure print_meaning;
%     begin print_cmd_chr(cur_cmd,cur_chr);
%     if cur_cmd>=call then
%       begin print_char(":"); print_ln; token_show(cur_chr); end
%     else if cur_cmd=top_bot_mark then
%       begin print_char(":"); print_ln; token_show(cur_mark[cur_chr]); end;
%     end;
%
% The `print_ln` is unconditional, so `\show\m` takes two lines — the name and
% `macro:` on the first, the parameter text and the body on the second. And
% `\meaning\m` takes one, because §470 builds `\meaning` with
% `selector=new_string` and §58's `print_ln` does nothing in that mode.
%
% The two share one string in this engine and both had the one-line form, so
% every `\show` of a macro was a line short. `print_cmd_chr` for a macro is
% `macro`, `\long macro` or `\protected\long macro` — never a colon — so the
% colon `print_char` writes is the first one in the string, and a body that
% contains a colon of its own stays where it belongs.
\catcode`\{=1 \catcode`\}=2 \catcode`\#=6
\tracingonline=1
\def\m{body}
\long\def\n#1{x#1y}
\def\c#1[#2]{a:b}
\let\z\m
\chardef\ch=65
\show\m \show\n \show\c  \show\z \show\ch \show\relax \show\topmark
% §1294 prints `sprint_cs(cur_cs); =` only when `cur_cs<>0`. Letters and
% other-character tokens therefore start directly with their meaning, while
% an active character still has a control-sequence name.
\show A \show!
\catcode`\~=13 \def~{active}\show~
\message{[\meaning\m][\meaning\n][\meaning\c]}
% A report last, so TeX's closing `)` — written against the current column —
% does not land on the message above. This engine writes no file echo at all.
\showthe\escapechar
\end
